\kbvcontinued ಯಾವುದೇ ವಾಕ್ಯದಲ್ಲಿ ಕರ್ತೃವು ಮಾಡುವ ಕೆಲಸವನ್ನು ತಿಳಿಸುವ ಫಲವನ್ನು ಕ್ರಿಯಾನಾಮವಾಗಿರುವುದು. ಮೇಲಿನ ಮೂರು ಮುಖಗಳಲ್ಲಿ ಕ್ರಮವಾಗಿ ಮಾಡಿದನು, ಹೊಲಿಯುತ್ತಾಳೆ ಬರೆದನು - ಕ್ರಿಯಾನಾಮಗಳಾಗಿರುವುವು. ಕರ್ತೃವಿನಿಂದ ಆಗುವ ಕಾರ್ಯ ಈ ಪದಗಳಲ್ಲಿ ವ್ಯಕ್ತವಾಗಿದೆ.\par

\vspace{2.5pt plus 3pt minus 0.5pt}\noindent\hspace*{1.2em}\kbvnum{(3)}\hspace{0.55em}ಕರ್ಮನಾಮ\par

\indent ಕತ್ತೆಯು ಭಾರವನ್ನು ಹೊರುತ್ತದೆ - - ಎಂಬ ವಾಕ್ಯಕ್ಕೆ ಸಂಬಂಧಿಸಿದಂತೆ\par

\indent ಕೇಳಬಹುದಾದ ಮೂರು ಪ್ರಶ್ನೆಗಳು ಇಲ್ಲಿವೆ. ಈ ಪ್ರಶ್ನೆಗಳನ್ನು ಗಮನಿಸಿ,\par

\indent ({\kbvlatin i}) ಭಾರವನ್ನು ಹೊರುವುದು ಯಾವುದು ? --- ಕ್ರಿಯಾನಾಮವನ್ನು ಅವಲಂಬಿಸಿ ಕೇಳಲಾಗುವ ಈ ಪ್ರಶ್ನೆಗೆ “ಕತ್ತೆಯು” ಎಂಬ ಉತ್ತರ ಸಹಜವಾಗಿ ಬರುವುದು. ಈ ರೀತಿಯಾಗಿ ಪ್ರಶ್ನಿಸಿ ಕರ್ತೃಪದವನ್ನು ಸುಲಭವಾಗಿ ಕಂಡುಹಿಡಿಯಬಹುದು.\par

\indent ({\kbvlatin ii}) ಕತ್ತೆಯು ಮಾಡುವ ಕಾರ್ಯವೇನು ? --- ಕರ್ತೃ ಮಾಡುವ ಕೆಲಸವೇನೆಂದು ಕೇಳುವ ಈ ಪ್ರಶ್ನೆಗೆ “ಹೊರುವುದು” ಎಂಬ ಉತ್ತರ ಸಹಜವಾಗಿ ಬರುವುದು. ಈ ರೀತಿಯಾಗಿ ಪ್ರಶ್ನಿಸಿ ಕ್ರಿಯಾನಾಮವನ್ನು ಸುಲಭವಾಗಿ ಕಂಡುಹಿಡಿಯಬಹುದು.\par

\indent ({\kbvlatin iii}) ಕತ್ತೆಯು ಏನನ್ನು ಹೊರುವುದು ? --- ಕ್ರಿಯಾಪದವನ್ನೇ ಅವಲಂಬಿಸಿ ಕೇಳಲಾಗುವ ಈ ಪ್ರಶ್ನೆಗೆ “ಭಾರವನ್ನು” ಎಂಬ ಉತ್ತರ ಸಹಜವಾಗಿ ಬರುವುದು. ಕ್ರಿಯಾತ್ಮಕ ಕರ್ತೃವಿನ ಇಚ್ಛೆ ಅಥವಾ ಇಚ್ಛಿತ ಕಾರ್ಯ ಯಾವುದೆಂಬುದು ಈ ಉತ್ತರದಿಂದ ತಿಳಿದುಬರುವುದು.\par

\indent ಹೀಗೆ ಕ್ರಿಯಾಪದವನ್ನು ಅವಲಂಬಿಸಿ ಏನನ್ನು? ಅಥವಾ ಯಾರನ್ನು? ಅಥವಾ ಯಾವುದನ್ನು? ಎಂಬ ಪ್ರಶ್ನೆ ಹಾಕಿದಾಗ ಉತ್ತರವಾಗಿ ಬರುವ ಪದವನ್ನು ಕರ್ಮನಾಮ ಎನ್ನುವರು.\par

\indent ಮೇಲೆ ತಿಳಿಸಿರುವ ವಾಕ್ಯಗಳಲ್ಲಿ “ಮೇಜನ್ನು”, “ಬಟ್ಟೆಗಳನ್ನು”, “ರಾಮಾಯಣವನ್ನು”... ಕರ್ಮನಾಮಗಳಾಗಿರುವುವು.\par
